\apendice{Identidades trigonometricas}

\subsection{Identidades fundamentales}

\begin{proposición}[Identidad pitagórica]
$$\sin^2(\alpha) + \cos^2(\alpha) = 1$$
\end{proposición}

\begin{proposición}[Identidades derivadas de la identidad pitagórica]
$$\tan^2(\alpha) + 1 = \sec^2(\alpha)$$
$$1 + \cot^2(\alpha) = \csc^2(\alpha)$$
\end{proposición}

\begin{proposición}[Definiciones de las funciones trigonométricas recíprocas]
$$\sec(\alpha) = \frac{1}{\cos(\alpha)}, \quad \csc(\alpha) = \frac{1}{\sin(\alpha)}, \quad \cot(\alpha) = \frac{1}{\tan(\alpha)}$$
\end{proposición}

\begin{proposición}[Relación tangente y cotangente]
$$\tan(\alpha) = \frac{\sin(\alpha)}{\cos(\alpha)}, \quad \cot(\alpha) = \frac{\cos(\alpha)}{\sin(\alpha)}$$
\end{proposición}

\subsection{Fórmulas de ángulo doble}

\begin{proposición}[Seno del ángulo doble]
$$\sin(2\alpha) = 2\sin(\alpha)\cos(\alpha)$$
\end{proposición}

\begin{proposición}[Coseno del ángulo doble]
$$\cos(2\alpha) = \cos^2(\alpha) - \sin^2(\alpha) = 2\cos^2(\alpha) - 1 = 1 - 2\sin^2(\alpha)$$
\end{proposición}

\begin{proposición}[Tangente del ángulo doble]
$$\tan(2\alpha) = \frac{2\tan(\alpha)}{1 - \tan^2(\alpha)}$$
\end{proposición}

\subsection{Fórmulas de ángulo triple}

\begin{proposición}[Seno del ángulo triple]
$$\sin(3\alpha) = 3\sin(\alpha) - 4\sin^3(\alpha)$$
\end{proposición}

\begin{proposición}[Coseno del ángulo triple]
$$\cos(3\alpha) = 4\cos^3(\alpha) - 3\cos(\alpha)$$
\end{proposición}

\begin{proposición}[Tangente del ángulo triple]
$$\tan(3\alpha) = \frac{3\tan(\alpha) - \tan^3(\alpha)}{1 - 3\tan^2(\alpha)}$$
\end{proposición}

\subsection{Fórmulas de ángulo mitad}

\begin{proposición}[Seno del ángulo mitad]
$$\sin\left(\frac{\alpha}{2}\right) = \pm\sqrt{\frac{1 - \cos(\alpha)}{2}}$$
\end{proposición}

\begin{proposición}[Coseno del ángulo mitad]
$$\cos\left(\frac{\alpha}{2}\right) = \pm\sqrt{\frac{1 + \cos(\alpha)}{2}}$$
\end{proposición}

\begin{proposición}[Tangente del ángulo mitad]
$$\tan\left(\frac{\alpha}{2}\right) = \frac{\sin(\alpha)}{1 + \cos(\alpha)} = \frac{1 - \cos(\alpha)}{\sin(\alpha)} = \pm\sqrt{\frac{1 - \cos(\alpha)}{1 + \cos(\alpha)}}$$
\end{proposición}

\subsection{Fórmulas de suma y diferencia}

\begin{proposición}[Seno de la suma]
$$\sin(\alpha + \beta) = \sin(\alpha)\cos(\beta) + \cos(\alpha)\sin(\beta)$$
\end{proposición}

\begin{proposición}[Seno de la diferencia]
$$\sin(\alpha - \beta) = \sin(\alpha)\cos(\beta) - \cos(\alpha)\sin(\beta)$$
\end{proposición}

\begin{proposición}[Coseno de la suma]
$$\cos(\alpha + \beta) = \cos(\alpha)\cos(\beta) - \sin(\alpha)\sin(\beta)$$
\end{proposición}

\begin{proposición}[Coseno de la diferencia]
$$\cos(\alpha - \beta) = \cos(\alpha)\cos(\beta) + \sin(\alpha)\sin(\beta)$$
\end{proposición}

\begin{proposición}[Tangente de la suma]
$$\tan(\alpha + \beta) = \frac{\tan(\alpha) + \tan(\beta)}{1 - \tan(\alpha)\tan(\beta)}$$
\end{proposición}

\begin{proposición}[Tangente de la diferencia]
$$\tan(\alpha - \beta) = \frac{\tan(\alpha) - \tan(\beta)}{1 + \tan(\alpha)\tan(\beta)}$$
\end{proposición}

\subsection{Fórmulas de producto a suma}

\begin{proposición}[Producto de senos]
$$\sin(\alpha)\sin(\beta) = \frac{1}{2}\left[\cos(\alpha - \beta) - \cos(\alpha + \beta)\right]$$
\end{proposición}

\begin{proposición}[Producto de cosenos]
$$\cos(\alpha)\cos(\beta) = \frac{1}{2}\left[\cos(\alpha - \beta) + \cos(\alpha + \beta)\right]$$
\end{proposición}

\begin{proposición}[Producto de seno y coseno]
$$\sin(\alpha)\cos(\beta) = \frac{1}{2}\left[\sin(\alpha + \beta) + \sin(\alpha - \beta)\right]$$
\end{proposición}

\begin{proposición}[Producto de coseno y seno]
$$\cos(\alpha)\sin(\beta) = \frac{1}{2}\left[\sin(\alpha + \beta) - \sin(\alpha - \beta)\right]$$
\end{proposición}

\subsection{Fórmulas de suma a producto}

\begin{proposición}[Suma de senos]
$$\sin(\alpha) + \sin(\beta) = 2\sin\left(\frac{\alpha + \beta}{2}\right)\cos\left(\frac{\alpha - \beta}{2}\right)$$
\end{proposición}

\begin{proposición}[Diferencia de senos]
$$\sin(\alpha) - \sin(\beta) = 2\cos\left(\frac{\alpha + \beta}{2}\right)\sin\left(\frac{\alpha - \beta}{2}\right)$$
\end{proposición}

\begin{proposición}[Suma de cosenos]
$$\cos(\alpha) + \cos(\beta) = 2\cos\left(\frac{\alpha + \beta}{2}\right)\cos\left(\frac{\alpha - \beta}{2}\right)$$
\end{proposición}

\begin{proposición}[Diferencia de cosenos]
$$\cos(\alpha) - \cos(\beta) = -2\sin\left(\frac{\alpha + \beta}{2}\right)\sin\left(\frac{\alpha - \beta}{2}\right)$$
\end{proposición}

\begin{proposición}[Suma de tangentes]
$$\tan(\alpha) + \tan(\beta) = \frac{\sin(\alpha + \beta)}{\cos(\alpha)\cos(\beta)}$$
\end{proposición}

\begin{proposición}[Diferencia de tangentes]
$$\tan(\alpha) - \tan(\beta) = \frac{\sin(\alpha - \beta)}{\cos(\alpha)\cos(\beta)}$$
\end{proposición}

\subsection{Potencias de funciones trigonométricas}

\begin{proposición}[Reducción de seno al cuadrado]
$$\sin^2(\alpha) = \frac{1 - \cos(2\alpha)}{2}$$
\end{proposición}

\begin{proposición}[Reducción de coseno al cuadrado]
$$\cos^2(\alpha) = \frac{1 + \cos(2\alpha)}{2}$$
\end{proposición}

\begin{proposición}[Reducción de tangente al cuadrado]
$$\tan^2(\alpha) = \frac{1 - \cos(2\alpha)}{1 + \cos(2\alpha)}$$
\end{proposición}

\begin{proposición}[Seno al cubo]
$$\sin^3(\alpha) = \frac{3\sin(\alpha) - \sin(3\alpha)}{4}$$
\end{proposición}

\begin{proposición}[Coseno al cubo]
$$\cos^3(\alpha) = \frac{3\cos(\alpha) + \cos(3\alpha)}{4}$$
\end{proposición}

\begin{proposición}[Seno a la cuarta potencia]
$$\sin^4(\alpha) = \frac{3 - 4\cos(2\alpha) + \cos(4\alpha)}{8}$$
\end{proposición}

\begin{proposición}[Coseno a la cuarta potencia]
$$\cos^4(\alpha) = \frac{3 + 4\cos(2\alpha) + \cos(4\alpha)}{8}$$
\end{proposición}

\subsection{Identidades de paridad}

\begin{proposición}[Paridad del seno]
El seno es una función impar:
$$\sin(-\alpha) = -\sin(\alpha)$$
\end{proposición}

\begin{proposición}[Paridad del coseno]
El coseno es una función par:
$$\cos(-\alpha) = \cos(\alpha)$$
\end{proposición}

\begin{proposición}[Paridad de la tangente]
La tangente es una función impar:
$$\tan(-\alpha) = -\tan(\alpha)$$
\end{proposición}

\begin{proposición}[Paridad de la cotangente]
La cotangente es una función impar:
$$\cot(-\alpha) = -\cot(\alpha)$$
\end{proposición}

\begin{proposición}[Paridad de la secante]
La secante es una función par:
$$\sec(-\alpha) = \sec(\alpha)$$
\end{proposición}

\begin{proposición}[Paridad de la cosecante]
La cosecante es una función impar:
$$\csc(-\alpha) = -\csc(\alpha)$$
\end{proposición}

\subsection{Complementariedad y suplementariedad}

\begin{proposición}[Cofunción del seno]
$$\sin\left(\frac{\pi}{2} - \alpha\right) = \cos(\alpha)$$
\end{proposición}

\begin{proposición}[Cofunción del coseno]
$$\cos\left(\frac{\pi}{2} - \alpha\right) = \sin(\alpha)$$
\end{proposición}

\begin{proposición}[Cofunción de la tangente]
$$\tan\left(\frac{\pi}{2} - \alpha\right) = \cot(\alpha)$$
\end{proposición}

\begin{proposición}[Cofunción de la cotangente]
$$\cot\left(\frac{\pi}{2} - \alpha\right) = \tan(\alpha)$$
\end{proposición}

\begin{proposición}[Cofunción de la secante]
$$\sec\left(\frac{\pi}{2} - \alpha\right) = \csc(\alpha)$$
\end{proposición}

\begin{proposición}[Cofunción de la cosecante]
$$\csc\left(\frac{\pi}{2} - \alpha\right) = \sec(\alpha)$$
\end{proposición}

\begin{proposición}[Seno suplementario]
$$\sin(\pi - \alpha) = \sin(\alpha)$$
\end{proposición}

\begin{proposición}[Coseno suplementario]
$$\cos(\pi - \alpha) = -\cos(\alpha)$$
\end{proposición}

\begin{proposición}[Tangente suplementaria]
$$\tan(\pi - \alpha) = -\tan(\alpha)$$
\end{proposición}

\subsection{Periodicidad}

\begin{proposición}[Periodicidad del seno]
$$\sin(\alpha + 2\pi k) = \sin(\alpha), \quad k \in \Z$$
\end{proposición}

\begin{proposición}[Periodicidad del coseno]
$$\cos(\alpha + 2\pi k) = \cos(\alpha), \quad k \in \Z$$
\end{proposición}

\begin{proposición}[Periodicidad de la tangente]
$$\tan(\alpha + \pi k) = \tan(\alpha), \quad k \in \Z$$
\end{proposición}

\begin{proposición}[Periodicidad de la cotangente]
$$\cot(\alpha + \pi k) = \cot(\alpha), \quad k \in \Z$$
\end{proposición}

\subsection{Fórmulas de linearización}

\begin{proposición}[Linearización de $\sin(\alpha)\sin(\beta)\sin(\gamma)$]
$$\sin(\alpha)\sin(\beta)\sin(\gamma) = \frac{1}{4}\left[\sin(\alpha + \beta - \gamma) + \sin(\beta + \gamma - \alpha) + \sin(\gamma + \alpha - \beta) - \sin(\alpha + \beta + \gamma)\right]$$
\end{proposición}

\begin{proposición}[Linearización de $\cos(\alpha)\cos(\beta)\cos(\gamma)$]
$$\cos(\alpha)\cos(\beta)\cos(\gamma) = \frac{1}{4}\left[\cos(\alpha + \beta - \gamma) + \cos(\beta + \gamma - \alpha) + \cos(\gamma + \alpha - \beta) + \cos(\alpha + \beta + \gamma)\right]$$
\end{proposición}

\subsection{Identidades de Weierstrass}

\begin{proposición}[Sustitución de Weierstrass para el seno]
Haciendo $t = \tan\left(\frac{\alpha}{2}\right)$:
$$\sin(\alpha) = \frac{2t}{1 + t^2}$$
\end{proposición}

\begin{proposición}[Sustitución de Weierstrass para el coseno]
Haciendo $t = \tan\left(\frac{\alpha}{2}\right)$:
$$\cos(\alpha) = \frac{1 - t^2}{1 + t^2}$$
\end{proposición}

\begin{proposición}[Sustitución de Weierstrass para la tangente]
Haciendo $t = \tan\left(\frac{\alpha}{2}\right)$:
$$\tan(\alpha) = \frac{2t}{1 - t^2}$$
\end{proposición}

\subsection{Identidades con ángulos notables}

\begin{proposición}[Valores para $\alpha = \pi$]
$$\sin(\pi) = 0, \quad \cos(\pi) = -1, \quad \tan(\pi) = 0$$
\end{proposición}

\begin{proposición}[Valores para $\alpha = \frac{\pi}{2}$]
$$\sin\left(\frac{\pi}{2}\right) = 1, \quad \cos\left(\frac{\pi}{2}\right) = 0, \quad \tan\left(\frac{\pi}{2}\right) \text{ no está definida}$$
\end{proposición}

\begin{proposición}[Valores para $\alpha = \frac{\pi}{3}$]
$$\sin\left(\frac{\pi}{3}\right) = \frac{\sqrt{3}}{2}, \quad \cos\left(\frac{\pi}{3}\right) = \frac{1}{2}, \quad \tan\left(\frac{\pi}{3}\right) = \sqrt{3}$$
\end{proposición}

\begin{proposición}[Valores para $\alpha = \frac{\pi}{4}$]
$$\sin\left(\frac{\pi}{4}\right) = \frac{\sqrt{2}}{2}, \quad \cos\left(\frac{\pi}{4}\right) = \frac{\sqrt{2}}{2}, \quad \tan\left(\frac{\pi}{4}\right) = 1$$
\end{proposición}

\begin{proposición}[Valores para $\alpha = \frac{\pi}{6}$]
$$\sin\left(\frac{\pi}{6}\right) = \frac{1}{2}, \quad \cos\left(\frac{\pi}{6}\right) = \frac{\sqrt{3}}{2}, \quad \tan\left(\frac{\pi}{6}\right) = \frac{\sqrt{3}}{3}$$
\end{proposición}

\subsection{Identidades adicionales}

\begin{proposición}[Identidad de Mollweide (Primera forma)]
$$\frac{a + b}{\cos\left(\frac{A - B}{2}\right)} = \frac{c}{\sin\left(\frac{C}{2}\right)}$$
donde $a, b, c$ son los lados de un triángulo y $A, B, C$ son los ángulos opuestos.
\end{proposición}

\begin{proposición}[Identidad de Mollweide (Segunda forma)]
$$\frac{a - b}{\sin\left(\frac{A - B}{2}\right)} = \frac{c}{\cos\left(\frac{C}{2}\right)}$$
\end{proposición}

\begin{proposición}[Fórmula de Euler]
$$e^{i\alpha} = \cos(\alpha) + i\sin(\alpha)$$
\end{proposición}

\begin{proposición}[Expresiones con la fórmula de Euler]
$$\cos(\alpha) = \frac{e^{i\alpha} + e^{-i\alpha}}{2}, \quad \sin(\alpha) = \frac{e^{i\alpha} - e^{-i\alpha}}{2i}$$
\end{proposición}

\begin{proposición}[Identidad de Ptolomeo]
Para un cuadrilátero cíclico con lados $a, b, c, d$ y diagonales $p, q$:
$$pq = ac + bd$$
\end{proposición}

\begin{proposición}[Suma de cuadrados de seno y coseno multiplicados]
$$\sin^2(\alpha) + \cos^2(\alpha + \beta) = \cos^2(\alpha) + \sin^2(\alpha + \beta)$$
\end{proposición}

